\section{Publishing and Querying}
\label{sec:publishing}

\term{python -m ui.api} (port 8000) starts one FastAPI application that hosts
the SPARQL endpoint, the Linked Data routes and a React interface (\term{ui/}),
all on the same graph; the earlier Gradio~\cite{abid2019gradio} interface is
still served by \term{python -m vidbpedia serve}. At start-up the server loads
the full N-Triples dump into an in-memory rdflib graph, loads
\term{inferred.nt} into a set so that every triple can be labelled as asserted
or inferred, builds a search index over labels and alternative names that
ignores Vietnamese diacritics, precomputes the class tree, and loads a second
graph without the inferred triples.

\subsection{Dereferenceable IRIs}
\label{sec:deref}

The server follows the pattern of DBpedia and of the W3C Cool URIs
note~\cite{sauermann2008cooluris}. Looking up a resource IRI never returns
content directly. It answers \emph{303 See Other} and redirects either to an
HTML page or to an RDF document, depending on the \term{Accept} header
(Table~\ref{tab:routes}). The HTML page is \term{/?resource=<name>}, which
opens the entity page \term{/entity/<name>}, so pages can be shared.

\begin{table}[t]
\centering
\caption{Linked Data routes (\term{web/linked\_data.py}). An unknown name returns 404.}
\label{tab:routes}
\small
\begin{tabularx}{\textwidth}{@{}l>{\raggedright\arraybackslash}X@{}}
\toprule
Route & Response \\
\midrule
\term{GET /resource/\{name\}} & 303 to \term{/?resource=\{name\}} for browsers, or to
  \term{/data/\{name\}.ttl}, \term{.nt}, \term{.jsonld} or \term{.rdf} when \term{Accept} asks for
  Turtle, N-Triples, JSON-LD or RDF/XML; sends \term{Vary: Accept} \\
\term{GET /page/\{name\}} & 303 to the HTML page, like \term{dbpedia.org/page/…} \\
\term{GET /data/\{name\}.\{ext\}} & every triple with the resource as subject, and up to 2{,}000 with it
  as object, in the requested format; sends \term{Access-Control-Allow-Origin: *} \\
\term{GET /ontology/\{term\}} & the definition of a \term{vio:} class or property, in Turtle \\
\bottomrule
\end{tabularx}
\end{table}

Listing~\ref{lst:curl} shows an exchange recorded on the running server. As
browsers do, clients must percent-encode the non-ASCII characters of the path.
In the response, the \term{dbo:*} types and the edge \term{vio:hasPlayer} are
inferred; \term{dbo:Animal} appears because the DBpedia Ontology places
\term{dbo:Person} under \term{dbo:Animal}, and our alignment module copies that
class chain.

\begin{listing}[t]
\begin{Verbatim}
$ curl -si -H "Accept: text/turtle" \
       http://127.0.0.1:8000/resource/Nguy%E1%BB%85n_C%C3%B4ng_Ph%C6%B0%E1%BB%A3ng
HTTP/1.1 303 See Other
vary: Accept
location: /data/Nguy%E1%BB%85n_C%C3%B4ng_Ph%C6%B0%E1%BB%A3ng.ttl

$ curl -sL -H "Accept: text/turtle" http://127.0.0.1:8000/resource/Nguy%E1%BB%85n_…
vres:Mito_HollyHock vio:hasPlayer vres:Nguyễn_Công_Phượng .
…
vres:Nguyễn_Công_Phượng a dbo:Animal, dbo:Athlete, dbo:Person, dbo:SoccerPlayer, … ;
    rdfs:label "Nguyễn Công Phượng"@en, "Nguyễn Công Phượng"@vi ;
    …
    vio:birthProvince vres:Nghệ_An ;
    vio:careerStation vres:Nguyễn_Công_Phượng__1, … ;
    …
    owl:sameAs dbr:Nguyễn_Công_Phượng, wd:Q18045362 ;
\end{Verbatim}
\caption{Dereferencing a resource IRI with \term{curl}. Response headers are
abridged and the Turtle is an excerpt.}
\label{lst:curl}
\end{listing}

\subsection{SPARQL Endpoint and Terminal Client}
\label{sec:endpoint}

The route \term{/sparql} implements the SPARQL~1.1
Protocol~\cite{feigenbaum2013protocol}, so that curl, SPARQLWrapper or YASGUI
can query the graph. It takes the query as a GET parameter, a POST form field
or an \term{application/sparql-query} body, with the 19 project prefixes
predeclared, and returns SPARQL JSON (the default), XML or CSV for
\term{SELECT} and \term{ASK}, and Turtle, N-Triples, JSON-LD or RDF/XML for
\term{CONSTRUCT} and \term{DESCRIBE}, chosen by a \term{format} parameter, as
on DBpedia, or by the \term{Accept} header. It is read-only, answers an invalid
query with 400 and the error, rejects \term{FROM} and \term{SERVICE}, which
would fetch other URLs, and runs queries in a worker thread. The
non-standard parameter \term{inference=false} restricts a query to asserted
triples (0 \term{dbo:Person} instances instead of 659).

The command \term{python -m vidbpedia query} runs the same queries from a
terminal, on the loaded dump or, with \term{-{}-endpoint}, through a running
endpoint (Listing~\ref{lst:sparql}). Both use one module, \term{kg/query.py},
which replaces two serializers of rdflib~7.6: its XML results writer drops the
literals 0 and false, and its text table loses the \term{ORDER BY}.

\begin{listing}[t]
\begin{Verbatim}
$ curl -H "Accept: text/csv" http://127.0.0.1:8000/sparql --data-urlencode "query=
    SELECT ?stadium ?seats WHERE { ?stadium a dbo:Stadium ; dbo:seatingCapacity ?seats }
    ORDER BY DESC(?seats) LIMIT 3"
stadium,seats
http://vi.dbpedia.org/resource/Sân_vận_động_Cần_Thơ,44400
http://vi.dbpedia.org/resource/Sân_vận_động_Quốc_gia_Mỹ_Đình,40192
http://vi.dbpedia.org/resource/Sân_vận_động_Thiên_Trường,30000
$ python -m vidbpedia query "SELECT ?province (COUNT(?p) AS ?players) WHERE {
    ?p vio:birthProvince ?province } GROUP BY ?province ORDER BY DESC(?players) LIMIT 3"
province       | players
---------------+--------
vres:Nghệ_An   | 58
vres:Hà_Nội    | 48
vres:Thanh_Hóa | 42
(3 dòng)
\end{Verbatim}
\caption{The SPARQL endpoint called with \term{curl} (CSV results) and the
terminal client (table). The first query uses only DBpedia terms, so every
answer comes from inferred triples.}
\label{lst:sparql}
\end{listing}

\subsection{Web Interface}
\label{sec:ui}

The interface has five pages (Figures~\ref{fig:ui}--\ref{fig:qa}). A top-bar switch hides
inferred triples in every graph.
\begin{itemize}
  \item \textbf{Overview} summarises the dataset, its inferred share, the
        links and the validation result.
  \item \textbf{Ontology} lists all 8{,}624 resources under their most specific
        \term{vio:} class, DBpedia superclasses after a $\sqsubseteq$ (a root shows
        its chain, \term{vio:Person} $\sqsubseteq$ \term{dbo:Person} $\sqsubseteq$
        \term{dbo:Animal}), and the OWL axioms with the triples each produced.
  \item \textbf{Entity} renders a page modelled on \term{dbpedia.org/page/…}:
        links to Wikipedia, DBpedia and Wikidata, the class tree with an
        asserted or inferred badge for each class, the \term{owl:sameAs} and
        provenance links, RDF downloads in four formats, the key facts, a career
        timeline, a draggable neighbourhood graph, a relation tree up to three
        hops deep and the property tables.
  \item \textbf{Question answering} (Section~\ref{sec:qa}).
  \item \textbf{SPARQL} is an editor with the 19 prefixes predeclared, nine
        example queries, a switch for inferred triples and JSON or CSV downloads.
\end{itemize}

\begin{figure}[t]
\centering
\includegraphics[width=\textwidth]{ui_entity.png}
\caption{Entity page of Nguyễn Công Phượng (abstract hidden): the asserted class and the named
graph, classes marked asserted (\emph{khai báo}) or inferred (\emph{suy luận}), and the
Linked Data card. Labels are in Vietnamese.}
\label{fig:ui}
\end{figure}

\begin{figure}[t]
\centering
\includegraphics[width=\textwidth]{ui_entity_graph.png}
\caption{Lower part of the same page: key facts, career timeline and the draggable
neighbourhood graph, where the \term{playedFor} edges are dashed because they are inferred.}
\label{fig:ui-graph}
\end{figure}

\subsection{Natural-Language Question Answering}
\label{sec:qa}

\begin{figure}[t]
\centering
\includegraphics[width=\textwidth]{ui_ask.png}
\caption{Question answering with \term{gpt-4o} for ``which players have played for Than Quảng
Ninh?'': the answer, the progress of the six steps, the generated query, the ontology terms
with their inferred share, and 39 rows with reasoning against 0 without (step~1 and the
checks of step~3 hidden). Labels are in Vietnamese.}
\label{fig:qa}
\end{figure}

A language model writes a SPARQL query from a prompt with a schema summary,
rules (never guess IRIs, use the \term{playedFor} shortcut, skip former
provinces) and six example queries. A query that fails or returns no rows goes
back once with the error, and a second prompt answers from at most 30 rows;
any OpenAI-compatible API works.

\paragraph{Entity linking.} Step~1 needs no model. The labels (Vietnamese and
English) and redirect titles of the 1{,}176 entities of the main classes are
folded to lower case without diacritics and stored in two hash tables: whole
names, also without a bracketed qualifier, and the last two to four words of
each name, since Vietnamese names start with the type or the family name
(``Nguyễn Quang Hải''). Spans of the
question are looked up longest first; a matched word is not reused, spans of
generic words are skipped, and a name shared by up to six entities keeps them all.
Linking takes about 1\,ms.

The page shows the answer above the six steps behind it:
(1)~entity linking, whose IRIs are added to the prompt; (2)~each generation attempt, with feedback naming IRIs absent
from the graph or of the wrong class; (3)~checks of syntax, read-only
form, \term{FROM}, \term{SERVICE}, \term{LIMIT} and linked IRIs, with the
inferred share of each ontology term; (4)~execution with and without inferred
triples; (5)~the answer; and (6)~an evidence graph (Figure~\ref{fig:qa}).
The examples follow one performance rule: rdflib applies a \term{FILTER} only after
joining the whole group, so the lookup of an entity by name is written as a
subquery that runs first. For the career of Nguyễn Công Phượng this takes the
query from 10.7\,s to 0.27\,s.
